\documentclass[10pt,a4paper]{amsart}
\usepackage[margin=19mm]{geometry}
\usepackage{amsmath,amssymb,mathtools}
\usepackage[scr=rsfs,cal=euler]{mathalfa}
\usepackage{xfrac}
\usepackage[unicode,hidelinks,bookmarks=false]{hyperref}
\newcommand{\ie}{i.e.\ }
\newcommand{\cf}{cf.\ }
\newcommand{\resp}{respectively\ }
\newcommand{\Subsection}{\subsection}
\providecommand{\nobreakdash}{\nobreak\hbox{-}}
\setlength{\emergencystretch}{2em}
\multlinegap0pt
\usepackage{stmaryrd}
\usepackage{amssymb}
\usepackage[all]{xy}
\usepackage{tikz}
\usepackage{mathtools}
\usepackage{stackengine}
\newtheorem{proposition}{Proposition}
\newtheorem{corollary}{Corollary}
\newtheorem{ex}{Example}
\def\l{\textup{l}}
\def\q{\textup{q}}
\def\r{\textup{r}}
\title{On lifting representations and actions on curves of the metacyclic groups $C_{p^s}\rtimes C_m$}
\author{Huy Dang, Adrian Vasiu}
\date{Source: arXiv:2609.03191v1 (CC BY 4.0)}
\begin{document}
\maketitle

\noindent\emph{Source and licence.} On lifting representations and actions on curves of the metacyclic groups $C_{p^s}\rtimes C_m$. Version arXiv:2609.03191v1 (CC BY 4.0), distributed by arXiv under the Creative Commons Attribution 4.0 International licence, http://creativecommons.org/licenses/by/4.0/, as recorded in the arXiv licence metadata for this exact version and shown on the abstract page https://arxiv.org/abs/2609.03191v1. Full licence notice: \texttt{licenses/arxiv-2609.03191-CC-BY-4.0.txt}.\par\smallskip

\noindent\textit{Editorial scope.} Counted unit: the article's proposition P8 (source0.tex lines 1507--1565). The article's complete original proof of it is delivered with it (source0.tex lines 1567--1620, one \texttt{proof} environment of 54 lines). No mathematics is written, repaired or re-derived: the statement and the proof are the article's own lines, in the article's own notation, and no retained line is substituted or re-flowed.  Retained uncounted as the source-local context the proof needs: lines 1115--1145; lines 1180--1198; lines 1362--1381; lines 1370--1372; lines 1412--1414; lines 1431--1433; lines 1436--1453; lines 1499--1501.  Nothing was omitted from any retained span, so the package carries no omission record.  No retained line contains a citation, so the wrapper carries no bibliography and no bibliography entry.  No retained line contains a figure environment, an image-inclusion command or drawing code, so no image is delivered and the package declares no asset dependency.  Editorial formatting only, in addition: an AMS \texttt{amsart} preamble that restates the article's own declarations and macros used by the retained lines, namely the environments \texttt{proposition}, \texttt{corollary}, \texttt{ex} on separate counters, together with the standard packages those lines need.  The retained environments print in this document as Proposition 1, Corollary 1, Example 1 in order of appearance, whereas the article prints the same items with the numbering of its own full document.  The complete native source of the article as distributed in the arXiv e-print for this exact version is bundled unchanged as \texttt{sources/209275/source0.tex}.
\begin{proposition}\label{P3}
We consider a cyclic $C_{p^s}$-extension of $k((t))$, its lower jumps $l_1<\cdots <l_s$, and its upper jumps $u_1<\cdots<u_s$. Then the following properties hold. 

\medskip
{\bf (1)} The lower and upper jumps determine uniquely each other by the recursive identities $l_1:=u_1$ and $l_i:=l_{i-1}+p^{i-1}(u_i-u_{i-1})$ for each $i\in\llbracket2,s\rrbracket$. 

\smallskip
{\bf (2)} For each $j\in\llbracket1,s\rrbracket $ we have $l_j=u_1+\sum_{i=1}^{j-1} p^i(u_{i+1}-u_i)$.

\smallskip
{\bf (3)} For each $j\in\llbracket1,s\rrbracket$ we have identities 
$$\sum_{i=1}^j \frac{l_i}{p^i}=\frac{\sum_{i=1}^j p^{i-1}u_i}{p^j}$$
and
$$p^{j-1}u_j:=l_j+(p-1)\sum_{i=1}^{j-1}l_ip^{j-i-1}.$$

{\bf (4)} For $b\in\llbracket0,p^s-1\rrbracket$, write
\[
b=\delta_{1,b}+\delta_{2,b}p+\cdots+\delta_{s,b}p^{s-1}
\]
with $\delta_{i,b}\in\llbracket0,p-1\rrbracket$ for each $i\in\llbracket1,s\rrbracket$ and define
\[
N_{s,b}:=\sum_{i=1}^s p^{s-i}\bigl[p-1+(p-1-\delta_{i,b})l_i\bigr].
\] Then for each $b\in \llbracket1,p^s-1\rrbracket$, if $j\in\llbracket0,s-1\rrbracket$ is the largest such that $p^j\mid b$, then we have identities 
\begin{equation*}
\begin{aligned}
N_{s,b-1}-N_{s,b}&=p^{s-j-1}\Big[l_{j+1}-(p-1)\sum_{i=1}^j p^{j+1-i}l_i\Big]\\
&=p^{s-j-1}[l_1+\sum_{i=1}^j p^i(u_{i+1}-pu_i)].
\end{aligned}
\end{equation*}
In particular, $N_{s,b-1}-N_{s,b}\ge p^{s-j-1}l_1\ge l_1>0$ and therefore we have $\frac{l_{j+1}}{p^{j+1}}-(p-1)\sum_{i=1}^j\frac{l_i}{p^i}\ge\frac{l_1}{p^{j+1}}>0$.
\end{proposition}

\begin{corollary}\label{C3}
Let $l_0:=0$. For a strictly increasing sequence $l_1<\cdots<l_s$ of positive integers, the following statements are equivalent.

\medskip
{\bf (1)} There exists a cyclic $C_{p^s}$-extension of $k((t))$ whose lower jumps are $l_1<\cdots<l_s$.

\smallskip
{\bf (2)} The following conditions hold.

\medskip\noindent
{\bf (2.a)} We have $p\nmid l_1$.

\smallskip\noindent
{\bf (2.b)} For each $j\in\llbracket1,s-1\rrbracket$ we have $p^j\mid l_{j+1}-l_j$ and
$$l_{j+1}-l_j\ge (p-1)\sum_{i=0}^{j-1} p^{j-i}(l_{i+1}-l_i).$$

\noindent
{\bf (2.c)} For each $j\in\llbracket1,s-1\rrbracket$ such that $l_{j+1}-l_j> (p-1)\sum_{i=0}^{j-1} p^{j-i}(l_{i+1}-l_i)$, the positive integer $l_1+\sum_{i=1}^{j}\frac{l_{i+1}-l_i}{p^i}$ is not divisible by $p$.
\end{corollary}

\begin{proposition}\label{P6}
Suppose that $X\rightarrow X/G$ is an HKG cover. Let
$$l_1<\cdots<l_s$$
be the lower jumps of the wild cyclic subgroup $H=C_{p^s}$ at the wild branch point. Then the following properties hold.

\medskip
\smallskip
{\bf (1) (Bleher--Chinburg--Kontogeorgis)} For each $b\in\llbracket0,p^s-1\rrbracket$, with the notation of Proposition \ref{P3}(4), we have
\begin{equation}\label{EQ51}
d_b=\left\lfloor \frac{N_{s,b}}{p^s}\right\rfloor=\left\lfloor \frac{\sum_{i=1}^s p^{s-i}[p-1+(p-1-\delta_{i,b})l_i]}{p^s}\right\rfloor.
\end{equation}

{\bf (2)} The set $\{x^idx|i\in\llbracket0,d_b-2\rrbracket\}$ is a $k$-basis of $H^0\bigl(\mathbb P_k^1,\Omega_{\mathbb P_k^1}(D_b)\bigr)$. In particular, if $d_b\le 1$, then this $k$-vector space is $\{0\}$.

\smallskip
{\bf (3)} For $b\in\llbracket0,p^s-1\rrbracket$, we have $d_b=0$ if and only if $b=p^s-1$.

\smallskip
{\bf (4)} The sequence $(d_b)_{b\in\llbracket0,p^s\rrbracket}$ is non-increasing.
\end{proposition}

\begin{equation}
d_b=\left\lfloor \frac{N_{s,b}}{p^s}\right\rfloor=\left\lfloor \frac{\sum_{i=1}^s p^{s-i}[p-1+(p-1-\delta_{i,b})l_i]}{p^s}\right\rfloor.
\end{equation}

\begin{equation}\label{EQ53}
\mu_{b,a}=\left\lfloor\frac{d_b-1}{m}\right\rfloor+\mu_b(a),
\end{equation}

\begin{equation}\label{EQ55}
\mu_X(a,b)=\mu_{b-1,a}-\mu_{b,a}.
\end{equation}

\begin{ex}\normalfont\label{EX8}
Suppose that $X\rightarrow X/G$ is an HKG cover. 

\medskip
{\bf (1)} As $d_{p^s-1}=0$ by Proposition \ref{P6}(3), we have $\mu_{p^s-1,a}=0$ for each $a\in C_m$ by Equation (\ref{EQ53}). Therefore we have $\mu_X(a,p^s)=\mu_{p^s-1,a}=0$ by Equation (\ref{EQ55}). Hence $\jmath_{V_X}=0$.

\smallskip
{\bf (2)} As we have $p^s-2=(p-2)+\sum_{i=1}^{s-1} (p-1)p^i$, Equation (\ref{EQ51}) gives
\[
d_{p^s-2}
=
\left\lfloor \frac{1}{p^s}\Big[p^{s-1}l_1+\Big(\sum_{l=1}^s p^{s-l}(p-1)\Big)\Big]\right\rfloor
=
\left\lfloor \frac{p^s-1+p^{s-1}l_1}{p^s}\right\rfloor
=\left\lfloor \frac{l_1}{p}\right\rfloor+1,
\]
where for the last identity we used that $p\nmid l_1$ by Corollary \ref{C3}(2.a). 
\end{ex}

\begin{equation}\label{EQ57}
\mu_X(a,b)=\nu_a\bigl(\delta_s(b-1)-1\bigr)-\nu_a\bigl(\delta_s(b)-1\bigr)
\end{equation}

\begin{proposition}[{\bf The Interval of Integers Principle}]
\label{P8}
Suppose that $X\rightarrow X/G$ is an HKG cover. For each integer $j\in\llbracket0,s-1\rrbracket$ let the pair $(\l_{j,0},\l_{j,-1})\in (\mathbb N\cup\{0\})\times\llbracket0,m-1\rrbracket$ be such that 
$$\Big\lfloor\frac{l_{j+1}}{p^{j+1}}-(p-1)\sum_{i=1}^j\frac{l_i}{p^i}\Big\rfloor=m\l_{j,0}+\l_{j,-1};$$
e.g., $(\l_{0,0},\l_{0,-1})\in (\mathbb N\cup\{0\})\times\llbracket0,m-1\rrbracket$ is such that $\lfloor\frac{l_1}{p}\rfloor=m\l_{0,0}+\l_{0,-1}$. Let $b\in\llbracket1,p^s\rrbracket$ and let the pair $(\q_{s,b},\r_{s,b})\in \mathbb Z\times\llbracket0,m-1\rrbracket$ be such that $\delta_s(b-1)-\delta_s(b)=m\q_{s,b}+\r_{s,b}$.  Then the following properties hold.

\medskip
{\bf (1)} If either $a=[0]_m$ with $\delta_s(b)\ge 1$ (i.e., with $b\le p^s-2$) or $a\neq [0]_m$, then we have identities
\begin{equation}\label{EQ58}
\begin{split}\mu_X(a,b)=&\q_{s,b}+\Big\lfloor\frac{\delta_s(b)-1+\r_{-a}+\r_{s,b}}{m}\Big\rfloor-\Big\lfloor\frac{\delta_s(b)-1+\r_{-a}}{m}\Big\rfloor\\
=&\begin{cases}\q_{s,b}+1\quad\;\textup{if}\quad \r_{\delta_s(b)-1+\r_{-a}}+\r_{s,b}\ge m\\
\q_{s,b}\quad\quad\quad\textup{if}\quad \r_{\delta_s(b)-1+\r_{-a}}+\r_{s,b}<m.\end{cases}
\end{split}
\end{equation}
In particular, we have $\mu_X(a,b)=\q_{s,b}+1$ if and only if $\r_{s,b}\neq 0$ (i.e., $\r_{s,b}\in\llbracket1,m-1\rrbracket$) and $a\in [\delta_s(b)]_m+\{[i]_m|i\in\llbracket0,\r_{s,b}-1\rrbracket\}$.

{\bf (2)} We have identities 
\begin{equation*}
\q_{s,p^s-1}=\Big\lfloor\frac{\lfloor \frac{l_1}{p}\rfloor+1}{m}\Big\rfloor=\begin{cases}\l_{0,0}+1\quad\textup{if}\quad\l_{0,-1}=m-1\\
\l_{0,0}\quad\quad\;\;\,\textup{if}\quad \l_{0,-1}\in\llbracket0,m-2\rrbracket,\end{cases}
\end{equation*} 
$\r_{s,p^s-1}=\Big\lfloor \frac{l_1}{p}\Big\rfloor+1-m\Big\lfloor\frac{\lfloor \frac{l_1}{p}\rfloor+1}{m}\Big\rfloor$, and
\begin{equation}\label{EQ59}
\begin{split}\mu_X([0]_m,p^s-1)=\q_{s,p^s-1}+\Big\lfloor\frac{\delta_s(p^s-1)-1+\r_{s,p^s-1}}{m}\Big\rfloor-1-\Big\lfloor\frac{\delta_s(p^s-1)-1}{m}\Big\rfloor\\
=\l_{0,0}=\begin{cases}\q_{s,p^s-1}\quad\quad\quad\textup{if}\quad \r_{s,p^s-1}\ge 1\\
\q_{s,p^s-1}-1\quad\;\textup{if}\quad \r_{s,p^s-1}=0\end{cases}=\begin{cases}\Big\lfloor\frac{\lfloor \frac{l_1}{p}\rfloor+1}{m}\Big\rfloor\quad\quad\quad\;\;\textup{if}\quad m\nmid \lfloor \frac{l_1}{p}\rfloor+1\\
\Big\lfloor\frac{\lfloor \frac{l_1}{p}\rfloor+1}{m}\Big\rfloor-1\quad\;\;\;\textup{if}\quad m\mid\lfloor \frac{l_1}{p}\rfloor+1.\end{cases}
\end{split}
\end{equation}

{\bf (3)} For each $a\in C_m$ we have inequalities 
$$\mu_X([0]_m,p^s-1)\le\mu_X(a,p^s-1)\le \mu_X([0]_m,p^s-1)+1.$$

{\bf (4)} For each pair $(a_1,a_2)\in C_m^2$ and every $b\in\llbracket1,p^s-1\rrbracket$ we have inequalities $0\le |\mu_X(a_1,b)-\mu_X(a_2,b)|\le 1$.

\smallskip
{\bf (5)} Let $j\in\llbracket0,s-1\rrbracket$. We assume that $b\in A_{s;j}$; so $p^j\mid b$ but $p^{j+1}\nmid b$. Then the following properties hold.

\medskip\noindent
{\bf (5.a)} We have identities
$$\q_{s,b}=\begin{cases}\l_{j,0}+1\quad\quad\textup{if}\quad\l_{j,-1}=m-1\quad\textup{and}\quad \r_{s,b}=0\\
\l_{j,0}\quad\quad\quad\;\;\,\textup{if}\quad\l_{j,-1}=m-1\quad\textup{and}\quad \r_{s,b}>0\\
\l_{j,0}\quad\quad\quad\;\;\,\textup{if}\quad \l_{j,-1}\in\llbracket0,m-2\rrbracket.\end{cases}$$

\noindent
{\bf (5.b)} If $j=0$ we assume that $b\neq p^s-1$. Then for each $a\in C_m$ we have inequalities 
$$\l_{j,0}\le \q_{s,b}\le\mu_X(a,b)\le \l_{j,0}+1\le\q_{s,b}+1.$$

\noindent
{\bf (5.c)} If $\l_{j,-1}=m-1$, $\r_{s,b}=0$, and for $j=0$ we have $b\neq p^s-1$, then $\mu_X(a,b)=\q_{s,b}=\l_{j,0}+1$.

\noindent
{\bf (5.d)} There exists a uniquely determined subset $\mathbb J_{X;j}$ of $C_m\times A_{s;j}$ such that 
$$\mu_X(a,b)=\begin{cases}\l_{j,0}+1\quad\;\textup{if}\quad (a,b)\in\mathbb J_{X;j}\\
\l_{j,0}\quad\quad\quad\textup{if}\quad (a,b)\in(C_m\times A_{s;j})\setminus\mathbb J_{X;j}.\end{cases}$$

{\bf (6)} Let $\mathbb J_X:=\sqcup_{j=0}^{s-1}\mathbb J_{X;j}$. We have an identity
$$\textup{g}_X=\frac{mp^s(p-1)}{2}\sum_{j=0}^{s-1} \l_{j,0}p^{s-1-j}+\sum_{(a,b)\in\mathbb J_X} b.$$
\end{proposition}

\begin{proof}
As $\delta_s(b-1)-1\ge 0$ by Proposition \ref{P6}(3), we have an identity $\nu_a\bigl(\delta_s(b-1)-1\bigr)=\lfloor\frac{\delta_s(b-1)-1+\r_{-a}}{m}\rfloor=\q_{s,b}+\lfloor\frac{\delta_s(b)-1+\r_{-a}+\r_{s,b}}{m}\rfloor$. 

For part (1) we have an identity $\nu_a\bigl(\delta_s(b)-1\bigr)=\lfloor\frac{\delta_s(b)-1+\r_{-a}}{m}\rfloor$ and hence part (1) follows from Equation (\ref{EQ57}).

For parts (2) to (5) we use the notation of Proposition \ref{P3}(4). 

We first prove part (5.a) when $j=0$. As $p\nmid b$, we have an identity $N_{s,b-1}=N_{s,b}+p^{s-1}l_1$ by Proposition \ref{P3}(4) and thus Equation (\ref{EQ51}) gives
\begin{equation}\label{EQ60}
\delta_s(b-1)-\delta_s(b)=\Big\lfloor\frac{l_1}{p}+\frac{N_{s,b}}{p^s}\Big\rfloor-\Big\lfloor\frac{N_{s,b}}{p^s}\Big\rfloor\in\Big\{\Big\lfloor\frac{l_1}{p}\Big\rfloor,\Big\lfloor\frac{l_1}{p}\Big\rfloor+1\Big\}.
\end{equation}
Thus $m\q_{s,b}+\r_{s,b}\in\{m\l_{0,0}+\l_{0,-1},m\l_{0,0}+\l_{0,-1}+1\}$ and from this we get that part (5.a) holds for $j=0$.

For part (2), we have $\delta_s(p^s-1)=0$ by Proposition \ref{P6}(3), $\r_{-[0]_m}=0$, and $\delta_s(p^s-2)=\left\lfloor \frac{l_1}{p}\right\rfloor+1$ by Example \ref{EX8}(2). Therefore $\q_{s,p^s-1}=\Big\lfloor\frac{\lfloor \frac{l_1}{p}\rfloor+1}{m}\Big\rfloor$ and $\r_{s,p^s-1}=\Big\lfloor \frac{l_1}{p}\Big\rfloor+1-m\Big\lfloor\frac{\lfloor \frac{l_1}{p}\rfloor+1}{m}\Big\rfloor$. The other identities of part (2) follow from them and the case $j=0$ of part (5.a). So part (2) holds. 

Based on part (1) and the case $j=0$ of part (5.a), to prove part (3) we can assume that $\mu_X([0]_m,p^s-1)=\q_{s,p^s-1}-1$. We have $m\mid\lfloor \frac{l_1}{p}\rfloor+1$, $\l_{0,-1}=m-1$, $\q_{s,p^s-1}=\l_{0,0}+1$, $\r_{s,p^s-1}=0$, and $\mu_X([0]_m,p^s-1)=\q_{s,p^s-1}-1=\l_{0,0}$ by part (2).  As $\r_{s,p^s-1}=0$, from part (1) we get that for each $a\in C_m\setminus\{[0]_m\}$ we have $\mu_X(a,p^s-1)=\q_{s,p^s-1}$. Thus part (3) holds.

Based on part (3), to prove part (4) we can assume that $b\le p^s-2$, and this case follows from part (1). So part (4) holds. 

Part (5.b) holds for $j=0$ by part (1). Thus to prove parts (5.a) and (5.b) we can assume that $s\ge 2$ and $j\in\llbracket1,s-1\rrbracket$. 

We have $N_{s,b-1}=N_{s,b}+l_{j+1}p^{s-j-1}-(p-1)\sum_{i=1}^j p^{s-i}l_i$ by Proposition \ref{P3}(4). Using Equation (\ref{EQ51}) we get that
$$\delta_s(b-1)-\delta_s(b)=\Big\lfloor \frac{l_{j+1}}{p^{j+1}}-(p-1)\sum_{i=1}^j\frac{l_i}{p^i}+\frac{N_{s,b}}{p^s}\Big\rfloor-\Big\lfloor\frac{N_{s,b}}{p^s}.\Big\rfloor$$ 
So $m\q_{s,b}+\r_{s,b}\in\{m\l_{j,0}+\l_{j,-1},m\l_{j,0}+\l_{j,-1}+1\}$. Thus, if $\l_{j,-1}\le m-2$, then $\q_{s,b}=\l_{j,0}$ and part (5.a) holds. 

If $\q_{s,b}=\l_{j,0}$, then from part (1) we get that part (5.b) holds. So to prove parts (5.a) and (5.b) we can assume that $\l_{j,-1}=m-1$. So we have the belonging relation $\delta_s(b-1)-\delta_s(b)\in\{m(\l_{j,0}+1)-1,m(\l_{j,0}+1)\}$.

We first assume that $\q_{s,b}>\l_{j,0}$. From the belonging relation and the inequalities $\delta_s(b-1)-\delta_s(b)\ge m\q_{s,b}\ge m(\l_{j,0}+1)$ we get the identity $\delta_s(b-1)-\delta_s(b)=m(\l_{j,0}+1)$. So $\q_{s,b}=\l_{j,0}+1$ and $\r_{s,b}=0$. As $\r_{s,b}=0$, we have $\mu_X=\q_{s,b}$ by part (1). So part (5.b) holds. 

Next we assume that $\q_{s,b}=\l_{j,0}$. From this and the belonging relation we get that $\delta_s(b-1)-\delta_s(b)=m(\l_{j,0}+1)-1$ and hence $m\nmid \delta_s(b-1)-\delta_s(b)$. 

From the last two paragraphs we get that part (5.a) holds if $\l_{j,-1}=m-1$. So part (5.a) holds. 

Part (5.c) follows from parts (5.a) and (5.b). 

Part (5.d) follows from part (5.b) and, when $b=p^s-1$, also from part (2). Thus part (5) holds.

For part (6), we first note that 
\begin{equation*}
\begin{aligned}
\sum_{b\in A_{s;j}} b&=p^j\Big(\sum_{i=1}^{p^{s-j}} i\Big)-p^{j+1}\Big(\sum_{i=1}^{p^{s-j-1}} i\Big)\\
&=\frac{p^{s}(p^{s-j}+1)}{2}-\frac{p^{s}(p^{s-j-1}+1)}{2}=\frac{p^{2s-j}-p^{2s-j-1}}{2}=\frac{p^{2s-j-1}(p-1)}{2}.\\
\end{aligned}
\end{equation*}
Based on this and $\mu_X(a,p^s)=0$ for each $a\in C_m$, we compute
\begin{equation*}
\begin{aligned}
\textup{g}_X&=\dim_k(V_X)=\sum_{(a,b)\in\mathbb J_{m,p^s}}\mu_X(a,b)b=\sum_{j=0}^{s-1} \sum_{(a,b)\in C_m\times A_{s;j}}\mu_X(a,b)b\\
&=\sum_{j=0}^{s-1} \Big[\sum_{(a,b)\in C_m\times A_{s;j}\setminus\mathbb J_{X;j}} \l_{j,0}b+\sum_{(a,b)\in \mathbb J_{X;j}} (\l_{j,0}+1)b\Big]\\
&=\Big(m\sum_{j=0}^{s-1} \l_{j,0}\sum_{b\in A_{s;j}} b\Big)+\sum_{j=0}^{s-1}\sum_{(a,b)\in \mathbb J_{X;j}}b\\
&=\frac{mp^s(p-1)}{2}\Big(\sum_{j=0}^{s-1} \l_{j,0}p^{s-1-j}\Big)+\sum_{(a,b)\in\mathbb J_X} b.\\
\end{aligned}
\end{equation*}
So part (6) holds.\end{proof}

\end{document}
